%!TEX root = ../Thesis.tex
\chapter{Intended Learning Objectives}

\noindent
The selected \acf{ILO} of this Thesis (also to be found in the study programme specification at \cite{GELOs}) can be found below, attached with statements of justified relevance, and listed in decreasing order of relevance to the project:

\begin{enumerate}

    \item has a thorough knowledge of basic mathematical and scientific methods which can be used to assess and solve idealized technical issues
    \note{The thesis will formalize image segmentation as a well-defined computational problem to be modelled mathematically through representations of images, filters, and segmentation masks, and evaluated quantitatively against ground truth data.}
    
    \item can examine, select and apply algorithms and data structures in a general engineering context
    \note{Existing state-of-the-art segmentation algorithms (SAM/MobileSAMv2) and traditional image processing methods will be examined, selected, and applied to a cross-disciplinary problem spanning digital humanities and computer vision.}
    
    \item can design and construct small and medium sized programs
    \note{The deliverable will be a structured proof-of-concept segmentation tool and pipeline, integrating pre-processing, semi-automatic segmentation, and export functionality.}

    \item is able to independently acquire new knowledge and adopt a critical approach to the acquired knowledge
    \note{Existing models and methods outside of the curriculum are identified, studied, and critically assessed with respect to their performance and suitability for the segmentation of Linear A inscriptions.}

    \item can combine research-based and practical knowledge to find suitable technological solutions, and see these in a societal context
    \note{Research-based computer vision methods will be applied to support the digitization workflow of SigLA, contributing to the public digital accessibility of the Linear A corpus.}
    
    \item has knowledge of the information structures of the subject and information sources relevant to the field of study, and can carry out relevant and critical information searches
    \note{A systematic literature review within computer vision and digital humanities will be conducted to situate the project within existing research.}

    \item is able to communicate technical information, theories, and results both graphically, in writing, and orally, and is able to present this to different groups of stakeholders
    \note{The computational pipeline will be documented, and segmentation results will be presented visually and analytically within the thesis.}

    \item can—on the basis of an independent professional approach—contribute to technical problem-solving through project work, independently as well as in collaboration with others
    \note{The evaluation of the tool will be made in close collaboration with the domain expert Ester Salgarella, while its development and documentation will be done alone.}


\end{enumerate}